% Gauss, Werke X/1 (1917), native leaves 33–34, printed 26–27.
\SourceLeaf{033}\NoSourceHead
\SourceGroup{II}
\HistoricTitle{[BESTIMMUNG DES DEZIDENTEN IN DER LEHRE VON DEN QUADRATISCHEN RESTEN.]}
\HistoricCitation{[Aus Handbuch 18, Bd. Mathematische Brouillons, October 1805, S. 164.]}
\begingroup
\newcommand{\Decident}[2]{\left(\frac{\underline{#1}}{#2}\right)}
\par\indent Designemus per $\Decident{k}{p}$ multitudinem productorum ex his
\[
k,\ 2k,\ 3k,\ \ldots,\ \tfrac12(p-1)k,
\]
\noindent quorum residua minima positiva secundum modulum $p$ cadunt infra $\tfrac12p$; porro per $[x]$ numerum integrum quantitate $x$ proxime minorem. Habebis
\begin{align*}
1)\qquad \Decident{k}{p}&=\left[\frac{2k}{p}\right]+\left[\frac{4k}{p}\right]+\left[\frac{6k}{p}\right]+\ldots+\left[\frac{(p-1)k}{p}\right]\\
&\quad-2\left[\frac{k}{p}\right]-2\left[\frac{2k}{p}\right]-2\left[\frac{3k}{p}\right]-\ldots-2\left[\frac{\tfrac12(p-1)k}{p}\right],\\[1ex]
2)\qquad [x]+[-x]&=-1,\\[1ex]
3)\qquad \Decident{k}{p}+\Decident{-k}{p}&=\tfrac12(p-1),\\[1ex]
4)\qquad \Decident{k}{p}&=\tfrac14(p\mp1)(k-1)-3\left[\frac{k}{p}\right]-\left[\frac{2k}{p}\right]-3\left[\frac{3k}{p}\right]-\ldots\\
&=\tfrac14(p\mp1)(k-1)\\
&\quad-2\left\{\left[\frac{k}{p}\right]+\left[\frac{3k}{p}\right]+\left[\frac{5k}{p}\right]+\ldots+\left[\frac{\tfrac12(p-3)k}{p}\right]\right\}\\
&\quad-\left\{\left[\frac{k}{p}\right]+\left[\frac{2k}{p}\right]+\left[\frac{3k}{p}\right]+\ldots+\left[\frac{\tfrac12(p-1)k}{p}\right]\right\}.
\end{align*}
\clearpage
\SourceLeaf{034}\ArithmeticHead{27}{R}{BESTIMMUNG DES DEZIDENTEN.}
\[
\begin{aligned}
\relax[x]+[2x]+[3x]+\ldots+[nx]&=0\left[\frac1x\right]\\
&\quad+1\left\{\left[\frac2x\right]-\left[\frac1x\right]\right\}\\
&\quad+2\left\{\left[\frac3x\right]-\left[\frac2x\right]\right\}\\
&\qquad\vdots\\
[nx]=h\qquad&\quad+h\left\{n-\left[\frac hx\right]\right\}
\end{aligned}
\]
\[
\left.\begin{aligned}
\relax[x]+[2x]+[3x]+\ldots+\overset{h}{[nx]}&\\
{}+\left[\frac1x\right]+\left[\frac2x\right]+\left[\frac3x\right]+\ldots+\left[\frac hx\right]&
\end{aligned}\right\}=nh.
\]
\begin{center}\textbf{BEMERKUNG.}\end{center}
\small
\par\indent Die vorstehende Aufzeichnung bildet eine Vorarbeit zu dem dritten Beweise des Reziprozitätsgesetzes (Comm. Gott. 16, 1808, Werke II, S. 1 ff.). Die hier bewiesene Formel bildet dort das \textit{Theorema} des art. 5 (Werke II, S. 7). Da sich auf folgenden Seiten des Handbuchs die Vorausberechnung der Örter von Juno und Ceres auf das Jahr 1808 findet, so wird diese Aufzeichnung wohl auf 1806—1807 anzusetzen sein. Vergl. den Artikel 20 des \textsc{Bachmann}schen Aufsatzes »Über \textsc{Gauss}' zahlentheoretische Arbeiten«.
\begin{flushright}\textsc{Schlesinger.}\end{flushright}
\ClosingRule
\vfill\hfill 4*\par
\clearpage\endgroup
